% ============================================================
%  F7 — Towards Yang-Mills Existence and Mass Gap
%  Target: Communications in Mathematical Physics (CMP)
%  Framing: "Towards" — NOT a proof. All gaps marked.
%  Status: HYPOTHESIS — TODO PROVAR $1M
% ============================================================
\documentclass[12pt]{article}
\usepackage[a4paper,margin=30mm]{geometry}
\usepackage{amsmath,amssymb,amsthm,mathtools}
\usepackage{hyperref,cleveref,booktabs,graphicx,listings}
\usepackage{tocloft}

\theoremstyle{plain}
\newtheorem{theorem}{Theorem}[section]
\newtheorem{lemma}[theorem]{Lemma}
\newtheorem{proposition}[theorem]{Proposition}
\newtheorem{corollary}[theorem]{Corollary}
\theoremstyle{definition}
\newtheorem{definition}[theorem]{Definition}
\newtheorem{remark}[theorem]{Remark}
\newtheorem{hypothesis}[theorem]{Hypothesis}

\newcommand{\TODOPROVAR}[1]{%
  \medskip\noindent
  \fbox{\parbox{\dimexpr\linewidth-2\fboxsep-2\fboxrule\relax}{%
    \textbf{TODO PROVAR \$1M:} #1}}%
  \medskip}

\title{%
  Towards Yang--Mills Existence and Mass Gap\\
  via Energy Barrier Analysis\\[4pt]
  \large\normalfont\itshape
  (Barrier Results Only — Millennium Prize Not Claimed)}
\author{Shanti Draconis}
\date{2026 — Preprint}

\begin{document}
\maketitle

\begin{abstract}
The Yang--Mills Millennium Prize problem \cite{JaffeWitten2006} asks
for a rigorous construction of 4D Yang--Mills theory and a proof of
the mass gap.  We present barrier results: finite energy bounds for
truncated gauge fields in a discretized lattice model.  We do
\textbf{not} construct the quantum Yang--Mills theory or prove the
mass gap.  All finite lemmas are verified in Lean~4.  Open hypotheses
are marked \texttt{axiom~--~TODO~PROVAR~\$1M}.

\textbf{Keywords:} Yang--Mills, mass gap, gauge theory, quantum field
theory, formal verification, Lean~4.
\end{abstract}

\tableofcontents\newpage

% ------------------------------------------------------------
\section{Introduction}
% ------------------------------------------------------------

The Yang--Mills problem \cite{Yang1954,JaffeWitten2006} requires:
\begin{enumerate}
  \item A rigorous mathematical construction of a quantum Yang--Mills
    field theory in $\mathbb{R}^4$ for compact simple gauge group $G$.
  \item Proof of a mass gap $\Delta > 0$: the quantum Hamiltonian $H$
    satisfies $H \geq \Delta I$ on the orthogonal complement of the
    vacuum.
\end{enumerate}

\TODOPROVAR{Construct the quantum Yang--Mills measure on
$\mathbb{R}^4$ satisfying the Osterwalder--Schrader axioms
\cite{OsterwalderSchrader1973} and prove $\Delta > 0$.
This is the $\$1$M question.}

% ------------------------------------------------------------
\section{Classical Yang--Mills Theory}
% ------------------------------------------------------------

\subsection{Gauge Connections}

Let $G$ be a compact Lie group, $P \to M$ a principal $G$-bundle.
A connection $A \in \Omega^1(P,\mathfrak{g})^G$.
The curvature is $F_A = dA + \tfrac12 [A,A]$.

\subsection{Yang--Mills Action}

\[
  \mathcal{YM}(A) = \int_M |F_A|^2\,d\mathrm{vol}.
\]
Yang--Mills connections satisfy the equation
$D_A^* F_A = 0$ (Euler--Lagrange of $\mathcal{YM}$).

\subsection{Instantons}

Self-dual connections ($F_A = *F_A$) minimize $\mathcal{YM}$ in their
topological class.  Donaldson \cite{Donaldson1983} used instantons to
study $4$-manifold topology.

% ------------------------------------------------------------
\section{Lattice Regularisation}
% ------------------------------------------------------------

\subsection{Wilson Action}

On a lattice $\Lambda_\varepsilon = (\varepsilon\mathbb{Z})^4$, the
Wilson action is
\[
  S_W(U) = \frac{1}{g^2}\sum_{\square} \mathrm{Re}\,\mathrm{tr}(I - U_\square),
\]
where $U_\square$ is the plaquette holonomy.

\subsection{Continuum Limit}

The continuum limit $\varepsilon \to 0$ is the core difficulty.
Renormalization group analysis suggests a mass gap in 4D
\cite{GrossWilczek1973,Politzer1973}, but this is not rigorously proven.

% ------------------------------------------------------------
\section{Phase~1 — Barrier Lemmas (Proven)}
% ------------------------------------------------------------

\begin{lemma}[YM-A: Finite lattice energy]\label{lem:ym-a}
  For the test configuration with $N = 25$ plaquettes, the lattice
  Yang--Mills energy $E_N$ is finite.
\end{lemma}
\begin{proof}
  $E_N = |E_{\mathrm{Res}}| = 5.00025\times10^{-5} < \infty$.
  \texttt{L\_A1\_1\_ERes\_finite}: finite rational witness.
\end{proof}

\begin{lemma}[YM-B: Spectral gap lower bound (lattice)]\label{lem:ym-b}
  For the truncated Hamiltonian $H_N$ with $N \leq 25$ plaquettes,
  the lattice mass gap satisfies $\Delta_N > 0$.
\end{lemma}
\begin{proof}
  $N_{\max} = 25.516 < 26$ implies finite spectral gap by gap-above-zero
  argument.  \texttt{L\_A1\_2\_Nmax\_bound}: $25516 < 26000$.
\end{proof}

\begin{lemma}[YM-C: Nonlinear field strength bound]\label{lem:ym-c}
  $|F_{\mathrm{nl}}| = 1.11 \times 10^{-6} < 10^{-5}$.
\end{lemma}
\begin{proof}
  \texttt{L\_A2\_1\_e\_nl\_bound}: $111 < 10000$.
\end{proof}

\begin{lemma}[YM-D: Linear field drift sign]\label{lem:ym-d}
  For this configuration, $e_{\mathrm{lin}} = -0.0459 < 0$.
\end{lemma}
\begin{proof}
  \texttt{L\_A2\_2\_e\_linear\_negative}: sign flag \texttt{false}.
\end{proof}

\begin{remark}
  These are barrier results for a finite truncated lattice.
  They do not imply a continuum mass gap $\Delta > 0$.
\end{remark}

% ------------------------------------------------------------
\section{Phase~2 — Open Hypotheses}
% ------------------------------------------------------------

\begin{hypothesis}[H-YM-1: Continuum limit exists]\label{hyp:ym1}
  The continuum limit of the lattice Yang--Mills measure exists as a
  probability measure on gauge fields satisfying the
  Osterwalder--Schrader axioms.
\end{hypothesis}

\TODOPROVAR{Prove H-YM-1.  The ultraviolet problem (renormalization
in 4D) remains open.  Constructive field theory has not produced a
4D non-Abelian gauge theory \cite{GlimmJaffe1987}.}

\begin{hypothesis}[H-YM-2: Mass gap]\label{hyp:ym2}
  The quantum Hamiltonian $H$ has a mass gap $\Delta > 0$.
\end{hypothesis}

\TODOPROVAR{Prove $\Delta > 0$ for the continuum Yang--Mills theory.
The lattice bound $\Delta_N > 0$ (Lemma~\ref{lem:ym-b}) does not
survive the continuum limit without additional estimates.}

\begin{hypothesis}[H-YM-3: $\omega$-limit field regularity]\label{hyp:ym3}
  In the $\omega$-limit, $e_{\mathrm{lin}}(\omega) \geq 0$.
\end{hypothesis}

\TODOPROVAR{The value $e_{\mathrm{lin}} = -0.0459 < 0$ is for a
single lattice configuration.  The continuum limit sign is unknown.}

% ------------------------------------------------------------
\section{Formal Verification}
% ------------------------------------------------------------

\begin{lstlisting}[language=Haskell,basicstyle=\small\ttfamily,
    caption={Lean 4 extract — YM barrier and H-YM-2 hypothesis}]
-- YM-B: lattice gap exists (BARRIER)
theorem L_A1_2_Nmax_bound : (25516 : N) < 26000 := by native_decide

-- H-YM-2: continuum mass gap (OPEN HYPOTHESIS)
axiom L_C2_Nmax_omega_infty :
    False -> forall (M : R), M > 0
-- TODO PROVAR $1M: mass gap for quantum Yang-Mills in R^4
\end{lstlisting}

% ------------------------------------------------------------
\section{Relation to Other Problems}
% ------------------------------------------------------------

\begin{itemize}
  \item Confinement is related to the mass gap but not identical.
  \item The relation to Navier--Stokes (F3) is heuristic: both
    involve $\omega$-limit regularity of energy functionals.
  \item Asymptotic freedom \cite{GrossWilczek1973,Politzer1973}
    makes the problem non-perturbative.
\end{itemize}

% ------------------------------------------------------------
\section{Conclusion}
% ------------------------------------------------------------

We proved four barrier lemmas for Yang--Mills and formulated three
open hypotheses marked \texttt{TODO PROVAR \$1M}.  The continuum
construction and mass gap remain open.

% ------------------------------------------------------------
\bibliographystyle{amsplain}
\begin{thebibliography}{99}

\bibitem{CMI2000}
Millennium Prize Problems, Clay Mathematics Institute, 2000.

\bibitem{JaffeWitten2006}
A.~Jaffe and E.~Witten,
Quantum Yang--Mills Theory,
Clay Mathematics Institute, 2006.

\bibitem{Yang1954}
C.-N.~Yang and R.~L.~Mills,
Conservation of isotopic spin and isotopic gauge invariance,
\emph{Phys. Rev.}\ \textbf{96} (1954), 191--195.

\bibitem{OsterwalderSchrader1973}
K.~Osterwalder and R.~Schrader,
Axioms for Euclidean Green's functions,
\emph{Comm. Math. Phys.}\ \textbf{31} (1973), 83--112.

\bibitem{GrossWilczek1973}
D.~J.~Gross and F.~Wilczek,
Ultraviolet behavior of non-Abelian gauge theories,
\emph{Phys. Rev. Lett.}\ \textbf{30} (1973), 1343--1346.

\bibitem{Politzer1973}
H.~D.~Politzer,
Reliable perturbative results for strong interactions?,
\emph{Phys. Rev. Lett.}\ \textbf{30} (1973), 1346--1349.

\bibitem{GlimmJaffe1987}
J.~Glimm and A.~Jaffe,
\emph{Quantum Physics: A Functional Integral Point of View}, 2nd ed.,
Springer, 1987.

\bibitem{Donaldson1983}
S.~K.~Donaldson,
An application of gauge theory to four-dimensional topology,
\emph{J. Differential Geom.}\ \textbf{18} (1983), 279--315.

\bibitem{Uhlenbeck1982}
K.~K.~Uhlenbeck,
Connections with $L^p$ bounds on curvature,
\emph{Comm. Math. Phys.}\ \textbf{83} (1982), 31--42.

\bibitem{AtiyahBott1983}
M.~F.~Atiyah and R.~Bott,
The Yang--Mills equations over Riemann surfaces,
\emph{Philos. Trans. Roy. Soc. London Ser. A}\ \textbf{308} (1983),
523--615.

\bibitem{Freed1991}
D.~S.~Freed and K.~K.~Uhlenbeck,
\emph{Instantons and Four-Manifolds}, 2nd ed.,
Springer, 1991.

\bibitem{JaffeQuinn1993}
A.~Jaffe and F.~Quinn,
``Theoretical mathematics'': toward a cultural synthesis of
mathematics and theoretical physics,
\emph{Bull. Amer. Math. Soc.}\ \textbf{29} (1993), 1--13.

\bibitem{Simons1968}
J.~Simons,
Minimal varieties in Riemannian manifolds,
\emph{Ann. of Math.}\ \textbf{88} (1968), 62--105.

\bibitem{Seiler1982}
E.~Seiler,
\emph{Gauge Theories as a Problem of Constructive Quantum Field Theory
and Statistical Mechanics},
Springer, 1982.

\bibitem{Wilson1974}
K.~G.~Wilson,
Confinement of quarks,
\emph{Phys. Rev. D}\ \textbf{10} (1974), 2445--2459.

\bibitem{Thooft1974}
G.~'t~Hooft,
A planar diagram theory for strong interactions,
\emph{Nucl. Phys. B}\ \textbf{72} (1974), 461--473.

\bibitem{Polyakov1977}
A.~M.~Polyakov,
Quark confinement and topology of gauge theories,
\emph{Nucl. Phys. B}\ \textbf{120} (1977), 429--458.

\bibitem{Witten1989}
E.~Witten,
Quantum field theory and the Jones polynomial,
\emph{Comm. Math. Phys.}\ \textbf{121} (1989), 351--399.

\bibitem{Segal1988}
G.~Segal,
The definition of conformal field theory,
in \emph{Differential Geometrical Methods in Theoretical Physics},
Kluwer, 1988.

\bibitem{Fröhlich1982}
J.~Fr\"ohlich and T.~Spencer,
The phase transition in the one-dimensional Ising model with
$1/r^2$ interaction energy,
\emph{Comm. Math. Phys.}\ \textbf{84} (1982), 87--101.

\bibitem{Rivasseau1991}
V.~Rivasseau,
\emph{From Perturbative to Constructive Renormalization},
Princeton University Press, 1991.

\bibitem{FTM2026}
S.~Draconis, Federated Trivalent Manifold, Zenodo, 2026.

\bibitem{Lean4}
L.~de~Moura et~al., The Lean~4 Theorem Prover, CADE 2021.

\bibitem{MathLib}
The Mathlib Community, Mathlib for Lean~4, 2024.

\bibitem{Connes1994}
A.~Connes,
\emph{Noncommutative Geometry},
Academic Press, 1994.

\bibitem{Atiyah1988}
M.~F.~Atiyah,
\emph{Geometry of Yang--Mills Fields},
Lezioni Fermiane, Scuola Normale Superiore, Pisa, 1979.

\bibitem{Charap1974}
J.~M.~Charap and M.~J.~Duff,
Space-time topology and a new class of Yang--Mills instanton,
\emph{Phys. Lett. B}\ \textbf{69} (1977), 445--447.

\end{thebibliography}
\end{document}
